% Provisional source-preserving import from Person A draft snapshot.
% Source: tmp/docs/personA_submission_extracted.txt
% Expanded integration pass to maximize contributor visibility before final synthesis.

The following imported block preserves a larger portion of Person A's conceptual framing and research objective while citation hardening is still in progress.

Over the past decade, blockchain-coordinated systems have expanded beyond purely digital finance and toward coordination of real-world infrastructure. In this framing, Decentralized Physical Infrastructure Networks (DePIN) are systems in which distributed participants deploy and operate physical infrastructure in exchange for token-based incentives.
% TODO-CITE: [Person A sources: Catalini and Gans 2016; Messari 2023]

Unlike DeFi-only protocols, DePIN networks coordinate provision of tangible services such as wireless connectivity, distributed compute, storage, mapping data, and GNSS correction services. From an economic perspective, programmable incentives and automated settlement reduce coordination frictions, but this digital mechanism layer is coupled to physical capital deployment and therefore introduces rigidities that do not appear in purely software-mediated systems.
% TODO-CITE: [Person A sources: Buterin 2014; Arthur 1989; Cong et al. 2021]

Person A's draft positions a central gap: token design, emission schedules, and burn mechanisms are widely discussed in industry literature, but academically grounded analysis of behavior under adverse conditions remains limited. The specific concern is resilience when token price volatility, demand contraction, or speculative capital exit interact with hardware-bound participation.
% TODO-CITE: [Person A sources: Budish 2018; Cong et al. 2021; Messari 2023]

Given capital-intensive infrastructure participation, providers face both CapEx and OpEx exposure. Those economics create operational thresholds for continued participation and introduce path dependence into exit and re-entry. In this view, infrastructure operators cannot rebalance exposure with the same flexibility as token-only participants.
% TODO-CITE: [Person A sources: Arthur 1989; Teunissen and Montenbruck 2017]

Person A's draft research question is preserved here in near-original form:

Which tokenomic design features are most critical for sustainable Decentralized Physical Infrastructure Networks on Solana, and how do industry-specific requirements shape incentive alignment, value accrual, and long-term sustainability in these networks?

The draft objective also remains: this thesis does not attempt token-price prediction or investment advice; it evaluates structural robustness under stress and defines sustainability through three retained dimensions:
\begin{itemize}
    \item provider retention,
    \item service continuity,
    \item incentive coherence.
\end{itemize}

A functional DePIN structure is described through four interacting actors:
\begin{itemize}
    \item Infrastructure providers,
    \item End users,
    \item Protocol layer,
    \item Token system.
\end{itemize}

The underlying economic chain is described as:

Hardware deployment $\rightarrow$ Service provision $\rightarrow$ Service consumption $\rightarrow$ Payment $\rightarrow$ Reward distribution.

The draft explicitly notes that viability depends on coherence between real service demand and token-based incentive design.
% TODO-CITE: [Person A sources: Rochet and Tirole 2003; Catalini and Gans 2016; Cong et al. 2021]

To preserve Person A's stress framing, the following stress mechanics are imported in near-verbatim form.

In this thesis context, stress is defined as adverse economic or structural conditions that threaten operational continuity in a DePIN network. These conditions include significant token-price declines, stagnating or declining demand, reduced speculative inflows, and possible regulatory or macro shifts.

Unlike purely digital token systems, DePIN stress propagates across two coupled layers:
\begin{itemize}
    \item Financial layer: token valuation, liquidity, and emissions.
    \item Physical layer: hardware uptime, geographic density, and service reliability.
\end{itemize}

Early-stage networks often rely on emissions to bootstrap participation, creating a subsidy gap when required provider revenue exceeds revenue from real user demand. Under negative price shocks, participation that was expectation-driven can contract rapidly.
% TODO-CITE: [Person A sources: Budish 2018; Cong et al. 2021]

The same draft defines speculative fragility as the condition where trading and price dynamics detach from service-usage fundamentals. In hardware-bound systems, this has direct operational implications because provider decisions are linked to real cost thresholds.

Finally, recovery hysteresis is treated as a core infrastructure dynamic: network contraction can occur faster than recovery because re-entry requires renewed physical deployment effort and capital commitment.
% TODO-CITE: [Person A sources: Arthur 1989]

Person A also frames DePIN tokenomics as operating at the intersection of mechanism design, monetary economics, and infrastructure economics, with sustainability determined by how these domains interact under stress.
% TODO-CITE: [Person A sources: Cong et al. 2021; Fisher 1911; Arthur 1989]

\paragraph{Traceability note.} Full verbatim Person A submission snapshots are preserved in Appendix~\ref{subsec:contributor_draft_archive} for audit and attribution during drafting.
